After these long terminological preliminaries, a little
theorem:

\begin{theorem}
\label{II.4.15}
Let $X$ be an $S$-prescheme locally of finite type, $Y$ a
closed subprescheme of~$X$ defined by a coherent sheaf
$\cal{J}$ of ideals on~$X$, $x$ a point of~$X$. One
now assumes \emph{$Y$ is smooth over~$S$ at $x$} (and nothing on~$X$). Then the following conditions are equivalent:
\begin{enumerate}
\item[(i)] $X$ is smooth over~$S$ at $x$
\item[(ii)] The immersion $i\colon Y \to X$ is regular at
$x$, \ie $\cal{J}_x$ is a regular ideal of~$\cal{O}_x$.
\end{enumerate}
\end{theorem}

\begin{corollary}
\label{II.4.16}
Suppose $Y$ is \emph{smooth} over~$S$. For $X$ to be smooth over~$S$
in a neighborhood of~$Y$ (\ie at the points of~$Y$) it is necessary and sufficient
that $Y$ be regularly immersed in $X$, \ie that
the immersion $i \colon Y \to X$ be regular.
\end{corollary}

\subsubsection*{Proof}
(i) implies (ii). One applies \ref{II.4.10} criterion~(ii), since
$g\colon X_1 \to X$ is \emph{flat}, to show that the inverse image
under $g$ of the subprescheme $Y'$ of~$X'$ is
regularly immersed, one is reduced to proving that
$Y'=S[t_{p+1},\dots, t_n]$ is regularly immersed in
$S[t_1,\dots,t_n]$, which is trivial (the $t_i$ ($1 \leq i \leq
p$) form a regular system of generators of
the Ideal defining $Y'$ in $X'$).

(ii) implies (i). Let $g_i$ ($1 \leq i \leq p$) be a regular
system of generators of~$\cal{J}_x$ and let $g_i$
($p+1 \leq i \leq n$) be elements of~$\cal{O}_{X,x}$ such that
their images $g'_i$ in $\cal{O}_{Y,x}$ define an
\emph{\'etale} morphism
\[
Y_1 \to Y'=S[t_{p+1},\dots, t_n]
\]
of a neighborhoods $Y_1$ of~$Y$ in $Y'$. The $g_i$ ($1 \leq i \leq n$)
come from sections (of the same name) of~$\cal{O}_X$ on a
neighborhood $X_1$ of~$x$, and one may suppose $X_1=X$, $Y_1=Y$. One
thus obtains a morphism
\[
g\colon X \to X'=S[t_1,\dots,t_n]
\]
and everything reduces to showing that this morphism is \emph{\'etale}
at~$x$. Taking $X_1$ small enough,
% original p. 48
one may suppose that the $g_i$ ($1 \leq i \leq p$) form a
regular system of generators of~$\cal{J}$ on all of
$X$. In particular, they generate $\cal{J}$, hence the
subprescheme $Y$ of~$X$ is identified with the inverse image
under $g$ of the subprescheme $Y'$ of~$X'$. Let $x'=g(x)$, then
the fiber of~$X' \to X$ at $x'$ is
identical to the fiber of~$Y \to Y'$ at $x$, hence is \'etale over~$\kres(x')$, hence $g$ is \emph{unramified} at $x$, it remains to
prove that $g$ is \emph{flat} at $x$. Now the associated graded
of $\cal{O}_{X',x'}$ filtered by the powers of
$\cal{J}'_{x}$ is \emph{free} over~$\cal{O}_{Y',x'}$ in all
degrees, on the other hand the associated graded of
$\cal{O}_{X,x}$ filtered by the powers of
$\cal{J}_x=\cal{J}'_x\cal{O}_{X,x}$ is isomorphic (under
the canonical homomorphism) to the tensor product of the former
with $\cal{O}_{Y,x}$ (since both rings are polynomial
rings in $n-p$ indeterminates, with ring of
constants $\cal{O}_{Y',x'}$, \resp $\cal{O}_{Y,x}$), finally over~$\cal{O}_{X',x'}/\cal{J}'_{x'}= \cal{O}_{Y',x'}$
$\cal{O}_{X,x}/\cal{J}_x =\cal{O}_{Y,x}$ is flat.

By a general flatness criterion (valid for
a local homomorphism of noetherian local rings $A' \to A$,
$A'$ being equipped with an ideal $J' \neq A'$ such that the associated
graded is free over~$A'/J'$ in every dimension) it follows that
$X$ is flat over~$X'$ at $x$, qed.

\begin{corollary}
\label{II.4.17}
Let $X$ be a prescheme locally of finite type over~$Y$, $i$
a section of~$X$ over~$Y$, $y$ a point of~$Y$, $x=i(y)$, $\cal{J}$
the sheaf of ideals on~$X$ defined by the
subprescheme $i(Y)$ (which we assume closed to
simplify the statement, a condition satisfied if $X$ is a
scheme).

The following conditions are equivalent:
\begin{enumerate}
\item[(i)] $X$ is smooth over~$Y$ at $x$
\item[(ii)] $i$ is a regular immersion at $y$
\item[(iii)] The completed $\cal{O}_y$-algebra of
$\cal{O}_x$ for the topology defined by the powers of
$\cal{J}_x$ is isomorphic to a formal power
series algebra $\cal{O}_y[[ t_1,\dots,t_n]]$.
\item[(iii~bis)] There exists an open neighborhood $U$ of~$y$ such that
the sheaf of algebras\linebreak
$\varprojlim i^*(\cal{O}_X/\cal{J}^{n+1})$
over~$\cal{O}_Y$ is isomorphic to a sheaf of the form
$\cal{O}_Y[[ t_1,\dots, t_n]]$ over~$U$.
\item[(iv)] There exists an open neighborhood $U$ of~$y$, and an
open neighborhood $V$ of~$x$, and finally a $Y$-morphism $g \colon V \to
U[t_1,\dots, t_n]$, such that $g$ is \'etale, that
$i$ induces a section of~$V$ over~$U$, transformed by $g$ into the
zero section of~$U[t_1,\dots, t_n]$ over~$U$.
\end{enumerate}
\end{corollary}

The equivalence
% original p. 49
of (i) and (ii) is a special case of
theorem~\ref{II.4.15}, setting $Y=S$, the equivalence of
(ii) and (iii) (and morally of (ii) and (iii~bis)) has been
pointed out with the ``recalls''. as for the equivalence of (i)
and (iv), it is easily deduced from
theorem~\ref{II.4.10} (equivalence of conditions (i) and
(ii) of the latter).

\begin{corollary}
\label{II.4.18}
Let $X$ be a prescheme smooth over~$S$. Then the
diagonal morphism
\[
\Delta_{X/S}\colon X\to X\times_S X
\]
is a
\emph{regular immersion},
or as one also says, $X$ is ``\emph{differentially
smooth}'' over~$S$.
\end{corollary}

Indeed, this is a special case of corollary~\ref{II.4.16}, since
$X$ and $X\times_S X$ are both smooth over~$S$.

% The corrected source repeats 4.18. Preserve it with a distinct PDF anchor.
\setcounter{theorem}{17}
\begingroup
\renewcommand{\theHtheorem}{\theHsection.\arabic{theorem}.bis}
\begin{remarks}
\label{rem:II.4.18}
Let us recall (I~\SourceRef{I.1}{1}) that if $X$ is a prescheme over
$S$, one introduces the quasi-coherent sheaves
of algebras
$\cal{P}_{X/S}^n=\cal{O}_{X\times_S X}/\cal{I}_X^{n+1}$ on~$X$,
(where $\cal{I}_X$ denotes the sheaf
of ideals
which defines the diagonal in $X\times_S X$), considered as
a sheaf of~$\cal{O}_X$-algebras thanks to the first
projection $\pr_1\colon X\times_S X\to X$. The
$\cal{P}_{X/S}^n$ form a projective system of Algebras over~$X$, whose projective limit is denoted $\cal{P}_{X/S}^{\infty}$
and is none other than the structure sheaf of the formal completion
of~$X\times_S X$ along the diagonal (now assuming $X$ is
locally of finite type over~$S$, hence the $\cal{P}_{X/S}^n$
are coherent). To say that $X$ is differentially smooth over~$S$
(\ie that the diagonal morphism $\Delta_{X/S}$ is a regular
immersion) also means that $\cal{P}_{X/S}^{\infty}$ is
regular, as a sheaf of algebras augmented to
$\cal{O}_X$, \ie that $\it{\Omega}_{X/S}^1$ is locally free and
the canonical surjective homomorphism
\[
S_{\cal{O}_X}(\it{\Omega}_{X/S}^1)\to
\mathrm{gr}_*(\cal{P}_{X/S}^{\infty})
\]
is an isomorphism, or finally that every point of~$X$
has an open neighborhood on which the sheaf of
augmented algebras $\cal{P}_{X/S}^{\infty}$ is isomorphic to a
sheaf $\cal{O}_X[[t_1,\dots,t_n]]$.

Let $s$ be a section of~$X$ over~$S$, $\cal{J}$ the sheaf
of ideals on~$X$ which it defines (assuming
% original p. 50
for simplicity that $s(S)$ is closed), one then has
canonical isomorphisms of augmented~$\cal{O}_X$-algebras:
\begin{equation}
\label{eq:II.4.4}
s^*(\cal{P}_{X/S}^n)=\cal{O}_X/\cal{J}^{n+1}\quoi,\quad
s^*(\cal{P}_{X/S}^{\infty})= \varprojlim_{n}
\cal{O}_X/\cal{J}^{n+1}
\end{equation}
These isomorphisms are functorial in an evident sense under
change of base, and (taking this fact into account) give back a
characterization of the sheaves of algebras $\cal{P}_{X/S}^n$ over~$S$. If for example $S=\Spec(k)$, $k$ a field, then the datum
of a section~$s$ of~$X$ over~$S$ is equivalent to the datum of a
point $x$ of~$X$ rational over~$k$, and the
preceding formulas mean that one has an isomorphism of
$k$-algebras
\begin{equation}
\label{eq:II.4.5}
\cal{P}_{X/S}^n(x)=\cal{O}_x/\goth{m}_x^{n+1}
\end{equation}
which justifies the name: ``\emph{sheaf of principal parts
of order~$n$ on~$X$ rel.\ to~$S$}''
given to~$\cal{P}_{X/S}^n$. One sees moreover from
\eqref{eq:II.4.4} that \emph{if $X$ is differentially smooth
over~$S$ at every point of~$s(S)$, then $X$ is smooth over~$S$ at every
point of~$s(S)$}, (corollary~\ref{II.4.17}) \emph{the converse
being equally true} (corollary~\ref{II.4.18}). Taking
\ref{II.4.13} into account, one easily concludes that if $X$ is an
$S$-prescheme locally of finite type, \emph{$X$ is smooth over~$S$ if and only if it is flat over~$S$ and differentially
smooth over~$S$.} (N.B. the flatness hypothesis is essential,
as one sees by taking for $X$ a closed subprescheme
of~$S$).

Let us note further, by way of recall, that one obtains a
\emph{second algebra structure} on~$\cal{P}_{X/S}^n$
thanks to the projection $\pr_2\colon X\times_S X\to X$,
which is moreover deduced from the former by means of
the \emph{canonical involution} of the sheaf of rings $\cal{P}_{X/S}^n$,
induced by the symmetry automorphism of~$X\times_S X$. One denotes
by $d_{X/S}^n$
or simply $d^n$, the homomorphism of sheaves of rings
\begin{equation}
\label{eq:II.4.6}
d_{X/S}^n\colon \cal{O}_X\to\cal{P}_{X/S}^n
\end{equation}
which corresponds to this second
Algebra structure. Taking the isomorphism~\eqref{eq:II.4.4} into account, this
homomorphism transforms a section $f$ of~$\cal{O}_X$ into a section
$d^n(f)$ of~$\cal{P}_{X/S}^n$ whose inverse image under a section $s$
of~$X$ over~$S$ is identified with the canonical image of~$f$ in
$\Gamma(X,\cal{O}_X/\cal{J}^{n+1})$. This justifies the name
``\emph{system of principal parts of order $n$ of~$f$}''
given to $d^nf$, notably in the case where $S=\Spec(k)$,
% original p. 51
envisaged in formula~\eqref{eq:II.4.5}.

To finish, let us note that the homomorphism \eqref{eq:II.4.6} may be
considered as the \emph{differential operator
of order}~$\leq n$\footnote{For everything concerning the present
paragraph, one may consult EGA~IV~16.8 to 16.12.}
(relative to the prescheme of constants $S$)
\emph{universal}
on~$\cal{O}_X$, by agreeing to call differential operator
of order~$\leq n$
from~$\cal{O}_X$ into a Module $F$, a homomorphism of sheaves $D$
which factors as
\[
D\colon \cal{O}_X\lto{d^n} \cal{P}_{X/S}^n\lto{u} F
\]
where $u$ is a homomorphism \emph{of~$\cal{O}_X$-Modules},
moreover uniquely determined by~$D$. This definition
agrees with the intuitive recursive definition ($D$ is a
differential operator of order $\leq n$ if for every section
$g$ of~$\cal{O}_X$ on an open set $U$ of~$X$, $f\mto D(fg)-D(f)$ is
a differential operator of order $\leq n-1$ on~$U$). It
follows that \emph{if $X$ is differentially smooth over~$S$, the
sheaf of rings of differential operators of all orders
has the simple well-known structure} in differential calculus on
differentiable manifolds, and in particular admits
locally a basis of~$\cal{O}_X$-Module
formed of the \emph{divided powers} in commuting operators
$\delta/\delta x_i$ $(1\leq i\leq n)$. If~$S$ is a sheaf of
$\QQ$-algebras ($\QQ=$ field of rationals) it suffices to take
the ordinary polynomials in the $\delta/\delta x_i$. In this case
moreover, and quite exceptionally, for $X$ to be
differentially smooth over~$S$, it already suffices that
$\it{\Omega}_{X/S}^1$ be locally free.
\end{remarks}
\endgroup

\begin{remark}
\label{II.4.19}
The terminology ``regular immersion'', ``regular
ideal'', etc.\ introduced in this no.\ has met with
a rather lively and general opposition (Chevalley, Serre). One has
proposed ``Cohen--Macaulay ideal'' or ``Macaulay
ideal'' or ``Macaulayan ideal'', which would morally also oblige
one to adopt ``Cohen--Macaulay immersion'' or ``Macaulay
immersion''. This terminology however conflicts with another
already used in future drafts of the
multiplodoque, where a morphism (of finite type) is said to be
``Cohen--Macaulay'' at a point if it is flat at that point, and if the
fiber passing through that point has there a local ring which is a
Cohen--Macaulay ring. Until a satisfactory solution is found,
we shall keep under all reservations the terminology introduced
in this no.\footnote{It is the one adopted in
EGA~$\mathrm{0}_\mathrm{IV}$~15.1.7.}
\end{remark}
